\documentclass[12pt]{article}

\usepackage{amsmath}
\usepackage{amsfonts}
\usepackage{graphicx}
\usepackage{float}
\usepackage{geometry}

\geometry{margin=1in}

\title{MNIST Classification using PCA, FDA, LDA and QDA}
\author{2024450}
\date{\today}

\begin{document}

\maketitle

\section{Introduction}

In this assignment we perform classification on a subset of the MNIST dataset. 
Only the handwritten digits \textbf{0, 1 and 2} are considered. 

The objective of this task is to:

\begin{itemize}
\item Reduce dimensionality using \textbf{Principal Component Analysis (PCA)}
\item Perform discriminative projection using \textbf{Fisher Discriminant Analysis (FDA)}
\item Compare classification performance of \textbf{Linear Discriminant Analysis (LDA)} and \textbf{Quadratic Discriminant Analysis (QDA)}
\item Analyze the impact of PCA on classification accuracy
\end{itemize}

Each MNIST image has a resolution of $28\times28$ pixels which corresponds to $784$ features. Working in such a high dimensional space is computationally expensive, therefore dimensionality reduction techniques are applied before classification.

\section{Dataset Description}

The MNIST dataset consists of handwritten digit images ranging from 0 to 9. 
In this experiment only the digits:

\[
\{0,1,2\}
\]

are used.

For each class:

\begin{itemize}
\item 100 samples are randomly selected for training
\item 100 samples are randomly selected for testing
\end{itemize}

Therefore,

\[
\text{Total Training Samples} = 300
\]

\[
\text{Total Test Samples} = 300
\]

Each image is converted into a vector of size:

\[
784 = 28 \times 28
\]

Thus the training data matrix becomes:

\[
X \in \mathbb{R}^{784 \times 300}
\]

\section{Data Preprocessing}

The following preprocessing steps were applied:

\begin{enumerate}

\item Images were converted to feature vectors by stacking pixels.
\item Pixel values were normalized to the range:

\[
[0,1]
\]

using

\[
x = \frac{x}{255}
\]

\item The dataset was organized into the matrix:

\[
X \in \mathbb{R}^{784 \times 300}
\]

\end{enumerate}

\section{Principal Component Analysis (PCA)}

PCA is used to reduce dimensionality while preserving the maximum variance of the data.

\subsection{Mean Calculation}

The mean vector of the dataset is computed as

\[
\mu = \frac{1}{N}\sum_{i=1}^{N}x_i
\]

where $N=300$.

\subsection{Mean Removal}

The centered data matrix is obtained as

\[
X_c = X - \mu
\]

\subsection{Covariance Matrix}

The covariance matrix is computed as

\[
S = \frac{X_c X_c^T}{N-1}
\]

\subsection{Eigen Decomposition}

The eigenvectors and eigenvalues of the covariance matrix are computed:

\[
S u_i = \lambda_i u_i
\]

The eigenvectors are sorted in descending order of eigenvalues.

\subsection{Dimensionality Reduction}

The projection matrix is constructed using the top $p$ eigenvectors:

\[
U_p = [u_1,u_2,\dots,u_p]
\]

The transformed data is:

\[
Y = U_p^T X_c
\]

The value of $p$ is chosen such that at least \textbf{75\% variance} is retained.

\section{Reconstruction}

The reconstructed samples are obtained by

\[
\hat{X} = U_p Y + \mu
\]

The reconstruction error is measured using Mean Squared Error (MSE):

\[
MSE = \frac{1}{d}\sum_{i=1}^{d}(x_i - \hat{x}_i)^2
\]

Reconstruction results for five samples are shown in output


\section{Fisher Discriminant Analysis (FDA)}

FDA is used to find a projection that maximizes class separability.

\subsection{Within Class Scatter}

\[
S_W = \sum_c \sum_{x_i \in C_c} (x_i - \mu_c)(x_i - \mu_c)^T
\]

\subsection{Between Class Scatter}

\[
S_B = \sum_c N_c(\mu_c - \mu)(\mu_c - \mu)^T
\]

where

\begin{itemize}
\item $\mu_c$ is the mean of class $c$
\item $N_c$ is the number of samples in class $c$
\end{itemize}

\subsection{Optimal Projection}

FDA finds projection matrix $W$ that maximizes:

\[
W = \arg\max \frac{\det(W^T S_B W)}{\det(W^T S_W W)}
\]

This leads to the generalized eigenvalue problem:

\[
S_W^{-1}S_B w = \lambda w
\]

The top eigenvectors form the projection matrix.

\section{Classification}

Two classifiers were evaluated.

\subsection{Linear Discriminant Analysis (LDA)}

The LDA discriminant function is:

\[
g_i(x) = x^T\Sigma^{-1}\mu_i - \frac{1}{2}\mu_i^T\Sigma^{-1}\mu_i
\]

where

\begin{itemize}
\item $\mu_i$ is the class mean
\item $\Sigma$ is the shared covariance matrix
\end{itemize}

The predicted class is

\[
\hat{y} = \arg\max_i g_i(x)
\]

\subsection{Quadratic Discriminant Analysis (QDA)}

QDA uses class-specific covariance matrices.

The discriminant function is:

\[
g_i(x) =
-\frac{1}{2}(x-\mu_i)^T\Sigma_i^{-1}(x-\mu_i)
-\frac{1}{2}\log|\Sigma_i|
\]

\section{Accuracy Metric}

Accuracy is defined as:

\[
Accuracy =
\frac{\text{Number of Correct Predictions}}
{\text{Total Number of Predictions}}
\]

Mathematically,

\[
Accuracy =
\frac{1}{N}\sum_{i=1}^{N} I(\hat{y_i} = y_i)
\]

where $I$ is the indicator function.


\section{Effect of PCA Variance}

Experiments were conducted with:

\begin{itemize}
\item 75\% variance
\item 90\% variance
\item first two principal components
\end{itemize}

Higher variance generally retains more information, leading to improved classification performance. However, increasing dimensionality also increases computational cost.

\section{Conclusion}

In this assignment we implemented PCA and FDA for dimensionality reduction and compared LDA and QDA classifiers.

Key observations:

\begin{itemize}

\item PCA successfully reduced the dimensionality of the dataset while retaining most of the variance.

\item FDA provided better class separation because it explicitly maximizes between-class variance while minimizing within-class variance.

\item LDA performed well when the covariance matrices of classes were similar.

\item QDA was more flexible but slightly more sensitive to noise due to class-specific covariance estimation.

\end{itemize}

Overall, combining PCA with discriminant analysis provides an effective pipeline for classification of high-dimensional image data.

\end{document}